\subsubsection{Annulations exactes de l’extracteur antipodal
}
\label{vi:subsubsec:cancelaciones-antipodales}

Les indices de cette sous-section se lisent dans \(\mathbb Z/12\mathbb Z\) et
on adopte \(\operatorname{sgn}(0)=0\). Soit \(A_r\) le décalage de
phase \((A_rT)(m)=T(m+r)\). Pour le registre neutre,

\begin{equation}
 e^\nu=(I-A_6)T_\nu,
 \qquad e^\nu_{m+6}=-e^\nu_m.
 \label{vi:eq:extractor-antipodal}
\end{equation}
On définit
\[
 B_a^\nu=e_a^\nu+e_{a+4}^\nu+e_{a+8}^\nu,
 \qquad
 \tau_m^\nu=\operatorname{sgn}(e_m^\nu),
 \qquad
 c_m^\nu=\tau_m^\nu\tau_{m+3}^\nu.
\]

\begin{proposition}[Annulation des deux lecteurs scalaires antipodaux
]
\label{vi:prop:cancelaciones-antipodales}
L’extracteur neutre satisfait

\begin{equation}
 B_{a+2}^\nu=-B_a^\nu,\qquad
 b_{120}^\nu:=\sum_{a=0}^{3}\operatorname{sgn}(B_a^\nu)=0,
 \label{vi:eq:cancelacion-b120}
\end{equation}
et
\begin{equation}
 c_{m+3}^\nu=-c_m^\nu,\qquad
 b_\zeta^\nu:=\sum_{m=0}^{11}c_m^\nu=0.
 \label{vi:eq:cancelacion-bzeta}
\end{equation}
\end{proposition}

\begin{proof}
Par~\eqref{vi:eq:extractor-antipodal},

\[
 B_{a+2}^\nu
 =e_{a+2}^\nu+e_{a+6}^\nu+e_{a+10}^\nu
 =-e_{a+8}^\nu-e_a^\nu-e_{a+4}^\nu=-B_a^\nu.
\]
Les quatre termes de \(b_{120}^\nu\) s’annulent dans les paires
\((0,2)\) et \((1,3)\), même lorsque l’un d’eux est nul. De même,
\(\tau_{m+6}^\nu=-\tau_m^\nu\), d’où

\[
 c_{m+3}^\nu
 =\tau_{m+3}^\nu\tau_{m+6}^\nu
 =-\tau_{m+3}^\nu\tau_m^\nu=-c_m^\nu.
\]
La somme des douze termes s’annule par paires alternées.

\end{proof}

Les annulations scalaires ne détruisent pas l’information orientée. Dans le
sous-espace anti-invariant

\[
 S_-:=\ker(A_6+I),
 \qquad J:=A_3\big|_{S_-},
\]
on a \(J^2=-I\). Comme le décalage est unitaire,
\(J^*=J^{-1}=-J\) : le registre conserve une structure complexe orientée.
De plus,

\[
 e^\nu=(I-A_6)T_\nu=2P_-T_\nu,
 \qquad P_-=(I-A_6)/2,
\]
de sorte que le vecteur antipodal complet reste disponible même si ses
deux contractions scalaires s’annulent.


Le retour cyclique \(A_{12}=I\) concerne cette représentation de phase ;
il n’équivaut pas au retour de l’état TPK complet, qui conserve retenue, feuillet,
frontière et mémoire. Pour la même raison,
\eqref{vi:eq:cancelacion-b120}--\eqref{vi:eq:cancelacion-bzeta} n’annulent pas
les coefficients généraux \(b_{120}\) et \(b_\zeta\) d’une route non
antipodale et ne modifient ni les masses, ni les signatures composées, ni le transport
PMNS. La couture longitudinale utilisée lors du prolongement du registre est
contrôlée indépendamment par le
lemme~\ref{vi:lem:costura-108-1080}.

